%% The explicit projection/substitution basis and restriction of models.

\section{Named presentations of Lawvere theories}

\begin{annexitem}{annex:p10.lawvere-present.1}{thm:p10-lawvere-presentation-recovery}{The equations are precisely the evaluation kernel}
\begin{proof}
\paragraph{Tuple identities and substitution.}
The selected power property gives
\[
 p_i^k\circ\langle g_0,\ldots,g_{k-1}\rangle=g_i,\qquad
                 \langle p_0^n,\ldots,p_{n-1}^n\rangle=1_n.
\]
It also gives
\[
 \langle a_0,\ldots,a_{r-1}\rangle\circ b
                 =\langle a_0\circ b,\ldots,a_{r-1}\circ b\rangle
\]
by comparing projections.  For an empty target list both sides
are the unique arrow to \(0\), and the empty tuple \(0\to0\)
is \(1_0\).

Let \(t\) be a term in \(X_k\), and let \(u_i\) be terms in
\(X_n\), for \(0\leq i<k\).  Put
\(a=\langle\llbracket u_0\rrbracket_n,\ldots,\llbracket u_{k-1}\rrbracket_n\rangle:n\to k\).
Structural induction proves the substitution identity
\[
          \bigl\llbracket t[u_0/x_0,\ldots,u_{k-1}/x_{k-1}]\bigr\rrbracket_n
                            =\llbracket t\rrbracket_k\circ a.
\]
For \(t=x_i\), this is the projection identity \(p_i^k a=\llbracket u_i\rrbracket_n\).
For \(t=\sigma_f(t_0,\ldots,t_{r-1})\), the induction hypotheses
and tuple precomposition give
\[
\begin{aligned}
 \bigl\llbracket t[\mathbf u/\mathbf x_k]\bigr\rrbracket_n
   &=f\circ\langle\llbracket t_0\rrbracket_k\circ a,\ldots,
                              \llbracket t_{r-1}\rrbracket_k\circ a\rangle\\
   &=f\circ\langle\llbracket t_0\rrbracket_k,\ldots,
                                  \llbracket t_{r-1}\rrbracket_k\rangle\circ a\\
   &=\llbracket t\rrbracket_k\circ a.
\end{aligned}
\]
For a nullary \(f:0\to1\), this calculation uses
\(!_k\circ a=!_n\).
For \(k=0\), substitution has no components and regards the
ground term in the new context; the same induction proves
\(\llbracket t\rrbracket_n=\llbracket t\rrbracket_0\circ!_n\).

\paragraph{Soundness of the equations.}
For any \(f:n\to1\),
\[
       \llbracket\sigma_f(\mathbf x_n)\rrbracket_n
             =f\circ\langle p_0^n,\ldots,p_{n-1}^n\rangle=f.
\]
This includes \(n=0\).
It proves soundness of each projection equation.
The two sides of a substitution equation both interpret as
\(f\circ\langle g_0,\ldots,g_{k-1}\rangle\).
In particular, at \(k=0\) both sides interpret as \(f\circ!_n\).

For each context \(X_n\), the relation of having equal
interpretations is an equivalence relation and is compatible
with every operation symbol.  The substitution identity shows
that it is preserved under all typed simultaneous substitutions.
Replacing substitution components by terms with the same
interpretations also preserves it, since their tuple arrows
are equal.  Thus these relations form a substitution-compatible
family of equational congruences containing \(E_{\mathcal T}\).
By the least-family definition,
\[
        s\equiv_{E_{\mathcal T},n}t
                        \ \Longrightarrow\
                     \llbracket s\rrbracket_n=\llbracket t\rrbracket_n.
\]
If a context has no terms, its component relation is simply empty;
the argument never requires that every term set be nonempty.

\paragraph{Reduction to a named operation.}
We prove by structural induction that
\[
             t\equiv_{E_{\mathcal T},n}
                           \sigma_{\llbracket t\rrbracket_n}(\mathbf x_n).
\]
For \(t=x_i\), use the projection equation in reverse.
Suppose \(t=\sigma_f(t_0,\ldots,t_{k-1})\), and put
\(g_i=\llbracket t_i\rrbracket_n\).  The induction hypotheses and
operation congruence give
\[
 t\equiv_{E_{\mathcal T},n}
       \sigma_f\bigl(\sigma_{g_0}(\mathbf x_n),\ldots,
                              \sigma_{g_{k-1}}(\mathbf x_n)\bigr).
\]
The defining substitution equation identifies the right side
with
\(\sigma_{f\circ\langle g_0,\ldots,g_{k-1}\rangle}(\mathbf x_n)\),
which is the required named operation.
If \(k=0\), there are no induction hypotheses; the nullary
instance of that same equation gives exactly the required step.

Consequently terms with the same interpretation reduce, by
equational derivation, to the same named term.  Symmetry and
transitivity yield the reverse implication to soundness.
This proves the displayed equality criterion, including \(n=0\).

\paragraph{The functor and its inverse.}
By soundness, the specified evaluation of a tuple is independent
of the representatives of its components.
Identity tuples evaluate to \(1_n\).
For composable term tuples \(U:n\to k\), \(V:k\to m\), the
substitution identity and tuple precomposition give
\[
          \varepsilon_{\mathcal T}(V\circ U)
                     =\varepsilon_{\mathcal T}(V)
                                      \circ\varepsilon_{\mathcal T}(U).
\]
If \(m=0\), all maps in this equation have terminal target and
their equality follows from uniqueness.  The case \(k=0\)
is included in the substitution identity.
Hence evaluation is an identity-on-objects functor.

It is faithful: equality of two evaluated tuples implies,
after composition with each target projection, equality of
the interpretations of corresponding components.  The
just-proved criterion equates these components in the term
theory.  For target \(0\) there is only the empty tuple.

For \(h:n\to m\), put \(h_j=p_j^m\circ h\).
The tuple of named terms
\((\sigma_{h_j}(\mathbf x_n))_{j<m}\) evaluates to
\[
            \langle h_0,\ldots,h_{m-1}\rangle
               =\langle p_0^m,\ldots,p_{m-1}^m\rangle\circ h=h.
\]
This proves fullness and gives the inverse displayed in the
statement.  When \(m=0\) it gives the unique empty tuple.
Applying this inverse after evaluation replaces every term
by its named operation, which is its own class by the
reduction already proved.  Thus the two maps are inverse on
every hom-set and on objects.

For completeness, the inverse also visibly respects composition.
If \(h:n\to k\), \(q:k\to m\), its \(j\)-th substituted
component is
\[
 \sigma_{p_j^m\circ q}\bigl(
       \sigma_{p_0^k\circ h}(\mathbf x_n),\ldots,
       \sigma_{p_{k-1}^k\circ h}(\mathbf x_n)\bigr).
\]
The substitution equation identifies this with
\(\sigma_{p_j^m\circ q\circ h}(\mathbf x_n)\), because
\(\langle p_i^k\circ h\rangle_{i<k}=h\).
The projection equations give the inverse functor's identity
law.  For \(n=0\) that identity tuple is empty.

\paragraph{Products, zero arities, and structural identifications.}
The term-theory projection \(n\to1\) is the class of the
one-component variable tuple \((x_i)\).
Its image is \(p_i^n\), since the one-fold product projection
\(p_0^1\) is \(1_1\).  Evaluation preserves all tuples by
definition, fixes \(0\), and consequently preserves the
selected finite products.  Its inverse preserves projections
by the projection equations.  A structural map \(n\to m\)
coming from a finite-ordinal function \(m\to n\) is a tuple
of the selected projections, so both functors preserve these
maps as well.

When \(n=0\) and \(h:0\to m\) exists with \(m>0\), each
\(p_j^m h:0\to1\) has its own nullary symbol, so the fullness
construction supplies ground terms for every component.
If \(\mathcal T(0,1)\) is empty, there are no nullary symbols
and hence no ground terms: a finite term tree must have a
leaf, but neither a variable nor a nullary symbol is available
in the empty context.  In this case there are also no arrows
\(0\to m\) for \(m>0\), since a target projection would
produce an arrow \(0\to1\).
For \(m=0\) there is always exactly one tuple and one arrow.

In a composition \(n\to0\to m\) with \(m>0\), the inverse
functor's second tuple consists of nullary terms.
Their substitution along \(n\to0\) is their use in \(X_n\).
The equation
\(\sigma_f()\approx\sigma_{f\circ!_n}(\mathbf x_n)\)
proves the composition equation componentwise.  This
explicitly covers passage through \(0\).

No step assumes distinct selected projections.
For example, if \(p_i^n=p_j^n\) in \(\mathcal T\), they name
the same operation symbol, and the two projection equations
give \(x_i\equiv_{E_{\mathcal T},n}x_j\).
The term theory therefore makes exactly the structural
identifications present in \(\mathcal T\).
\end{proof}
\end{annexitem}

\begin{annexitem}{annex:p10.lawvere-present.2}{thm:p10-lawvere-model-restriction}{Morphisms, named equations, and restriction of models}
\begin{proof}
Let \(H:\mathcal T\to\mathcal U\) be identity on objects and
preserve all the selected projections.
For any tuple \(a=\langle f_0,\ldots,f_{m-1}\rangle:n\to m\),
\[
 p_j^{m,\mathcal U}\circ H(a)
     =H(p_j^{m,\mathcal T}\circ a)=H(f_j).
\]
The selected product property in \(\mathcal U\) therefore gives
\[
                     H(\langle f_j\rangle_{j<m})
                                  =\langle H(f_j)\rangle_{j<m}.
\]
For \(m=0\), this is terminal uniqueness.  Hence \(H\)
preserves the selected powers and their tuples.
Products of any finite list of objects are obtained by
concatenating their finite lists of copies of \(1\);
their block projections are tuples of selected projections.
These products are therefore preserved too.
Any other finite product cone is isomorphic, compatibly
with its legs, to such a selected cone, and a functor sends
that isomorphism to an isomorphism.  Thus \(H\) preserves
finite products.

A structural finite-set map is a tuple of projections, so
preserving objects and projections implies commutation with
the structural functors.  Conversely, that commutation
implies the two requirements by applying it to objects and
the individual selectors.  This proves the equivalence in
Definition~\ref{def:p10-lawvere-theory-morphism}.
Identity functors and composites retain the requirements.

If \(M:\mathcal U\to\mathsf{Set}\) preserves finite products,
so does \(M\circ H\).  For a natural transformation
\(\alpha:M\Rightarrow N\), use the same component
\(\alpha_n:M(n)\to N(n)\) at each object of \(\mathcal T\).
Its naturality at \(a:n\to m\) is precisely the original
naturality equation at \(H(a)\):
\[
                    N(H(a))\circ\alpha_n
                              =\alpha_m\circ M(H(a)).
\]
Componentwise identities and composition show that this
defines \(H^*\) on the full categories of models.
Since \(H(1)=1\), its underlying carrier is literally \(M(1)\);
nonemptiness is preserved, and the restriction to the stated
full subcategories follows.
Associativity of functor composition and the unchanged
components give
\((K\circ H)^*=H^*\circ K^*\) and \((1_{\mathcal T})^*=1\).

Finally define a translation of named terms by fixing each
variable and recursively replacing
\(\sigma_f(t_0,\ldots,t_{k-1})\) by
\(\sigma_{H(f)}(\widehat Ht_0,\ldots,\widehat Ht_{k-1})\).
Identity on objects makes this an arity-preserving translation.
An image of a projection equation is a projection equation
for \(\mathcal U\).
The image of a substitution equation is one for \(\mathcal U\),
because
\[
 H\bigl(f\circ\langle g_0,\ldots,g_{k-1}\rangle\bigr)
           =H(f)\circ\langle H(g_0),\ldots,H(g_{k-1})\rangle.
\]
This includes empty tuples.
Translation commutes with simultaneous substitution by
induction on terms, and preserves every congruence operation.
It therefore descends to the quotient term theories and
preserves identities, composition, and selected projections.

Induction on a term also gives
\[
                  \llbracket\widehat Ht\rrbracket_n^{\mathcal U}
                                 =H(\llbracket t\rrbracket_n^{\mathcal T}):
\]
the variable step uses the projection condition, and the
operation step uses preservation of tuples and composition.
Taking tuples proves
\(\varepsilon_{\mathcal U}\circ\widehat H
       =H\circ\varepsilon_{\mathcal T}\) on all arrows and objects.
\end{proof}
\end{annexitem}
